\exercisesection{}

\begin{enumerate}
\def\labelenumi{\arabic{enumi})}
\tightlist
\item
  In \(\mathbf{R}^2\) , let \(\lambda\) and \(\mu\) be the positive measures
\end{enumerate}

\[
f \mapsto \int_ {0} ^ {+ \infty} f (x, 0) d x, \quad f \mapsto \int_ {0} ^ {+ \infty} f (- x, x) d x \quad \left(f \in \mathcal {K} (\mathbf {R} ^ {2})\right).
\]

Show that \(\lambda\) and \(\mu\) are convolvable. Let \(u\) be the homomorphism \((x,y) \mapsto x\) of \(\mathbf{R}^2\) onto \(\mathbf{R}\) . Show that \(u\) is \(\lambda\) -proper and \(\mu\) -proper, but that \(u(\lambda)\) and \(u(\mu)\) are not convolvable.

\begin{enumerate}
\def\labelenumi{\arabic{enumi})}
\setcounter{enumi}{1}
\item
  Let X be a locally compact space in which a locally compact group G operates on the left continuously. Let E be a linear subspace of \(\mathcal{M}(X)\) , stable for the \(\boldsymbol{\gamma}(s)\) ( \(s \in G\) ), equipped with a quasi-complete locally convex topology finer than the topology of compact convergence in \(\mathcal{K}(X)\) . For every \(s \in G\) , let \(\boldsymbol{\gamma}_{\mathrm{E}}(s)\) be the restriction of \(\boldsymbol{\gamma}(s)\) to E. Assume that \(\mu \in E\) implies \(|\mu| \in E\) , and that the representation \(\gamma_{E}\) of G in E is equicontinuous. Then, if \(\mu \in \mathbf{E}\) and \(\nu \in \mathcal{M}^1 (\mathbf{G})\) , \(\nu\) and \(\mu\) are convolvable and \(\nu * \mu = \gamma_{\mathbf{E}}(\nu)\mu \in \mathbf{E}\) . (Make use notably of Prop. 17 of Ch. VI, §1, No.~7.)
\item
  For \(x = (x_{1}, \ldots, x_{n}) \in \mathbf{R}^{n}\) one sets \(|x|^{2} = x_{1}^{2} + \cdots + x_{n}^{2}\) . Let \(M_{1}\) be the set of measures \(\mu\) on \(R^{n}\) for which there exists a real number k such that the function \((1 + |x|^{2})^{k}\) is \(\mu\) -integrable. Let \(M_{2}\) be the set of measures \(\nu\) on \(R^{n}\) for which the function \((1 + |x|^{2})^{k}\) is \(\nu\) -integrable for every k. Show that if \(\mu \in M_{1}\) and \(\nu \in M_{2}\) , then \(\mu\) and \(\nu\) are convolvable, and \(\mu * \nu \in M_{1}\) ; if \(\mu \in M_{2}\) and \(\nu \in M_{2}\) , then \(\mu * \nu \in M_{2}\) . (One first shows that if \(u \geqslant 0\) , \(v \geqslant 0\) , then
\end{enumerate}

\[
(1 + u ^ {2}) (1 + v ^ {2}) \geqslant \frac {1}{3} (1 + (u + v) ^ {2});
\]

from this, one deduces that for any \(x, y\) in \(\mathbf{R}^n\) ,

\[
1 + | x | ^ {2} \leqslant 3 (1 + | y | ^ {2}) (1 + | x + y | ^ {2}).
\]

Then let \(\mu \in \mathcal{M}_1\) , \(\nu \in \mathcal{M}_2\) with \(\mu \geqslant 0\) , \(\nu \geqslant 0\) . Let \(f\) be a continuous function \(\geqslant 0\) on \(\mathbf{R}^n\) . There exists a \(k\) such that \(\mu = (1 + |x|^2)^k \cdot \mu_1\) with \(\mu_1\) bounded; let \(\nu_1 = (1 + |x|^2)^k \cdot \nu\) , which is bounded. Then

\[
\int^ {*} f (x + y) d \mu (x) d \nu (y) \leqslant 3 ^ {k} \int^ {*} (1 + | x | ^ {2}) ^ {k} f (x) d (\mu_ {1} * \nu_ {1}) (x).
\]

Whence the convolvability of \(\mu\) and \(\nu\) and the fact that \(\mu * \nu \in \mathcal{M}_1\) . Argue in an analogous manner for \(\mu, \nu\) in \(\mathcal{M}_2\) .)

\begin{enumerate}
\def\labelenumi{\arabic{enumi})}
\setcounter{enumi}{3}
\item
  Let \(\mu\) be Lebesgue measure on \(\mathbf{R}\) , and \(\nu\) Lebesgue measure on \([0, +\infty[\) . Let \(x_1, x_2\) be in \(\mathbf{R}\) . Show that the convolution product \(((\varepsilon_{x_1} - \varepsilon_{x_2}) * \nu) * \mu\) is defined, but that \(\mu\) and \(\nu\) are not convolvable. Show that the convolution products \(\nu * ((\varepsilon_{x_1} - \varepsilon_{x_2}) * \mu)\) and \((\nu * (\varepsilon_{x_1} - \varepsilon_{x_2})) * \mu\) are defined, but are distinct for \(x_1 \neq x_2\) .
\item
  Let G be a compact group, and \(\mu\) a positive measure on G, with support G, such that \(\mu*\mu=\mu\) . Show that \(\mu\) is the normalized Haar measure of G. (First show that \(\|\mu\|=1\) . Then, if \(\mu\) is not the Haar measure of G, there exists a function \(f\in\mathcal{K}_{+}(G)\) such that
\end{enumerate}

\[
\int f (s) d \mu (s) \geqslant \int f (s t) d \mu (s)
\]

for all \(t \in \mathbf{G}\) , and such that \(\int f(s) d\mu(s) > \int f(st) d\mu(s)\) for some \(t\) . Show that one then has \((\mu * \mu)(f) > \mu(f)\) .

\begin{enumerate}
\def\labelenumi{\arabic{enumi})}
\setcounter{enumi}{5}
\tightlist
\item
  \begin{enumerate}
  \def\labelenumii{\alph{enumii})}
  \tightlist
  \item
    Let \(I = [0,1]\) . Let \(f\) be a continuous function \(\geqslant 0\) on \(\mathbf{R}\) , with support contained in \([-1,0]\) . Show that the set of functions \(\pmb{\gamma}(s)f|I\) ( \(s \in I\) ) has infinite rank in \(\mathcal{K}(I)\) .
  \end{enumerate}
\end{enumerate}

\begin{enumerate}
\def\labelenumi{\alph{enumi})}
\setcounter{enumi}{1}
\item
  Let \(f_1, \ldots, f_n\) be in \(\mathcal{K}(\mathbf{R})\) . Let \(M\) be the set of measures \(\mu \in \mathcal{M}(I)\) such that \(\mu(f_1) = \ldots = \mu(f_n) = 0\) . Show that the set of functions \(y \mapsto \int f(x + y) d\mu(x)\) ( \(y \in I\) ), where \(\mu\) runs over \(M\) , has infinite rank in \(\mathcal{K}(I)\) . (Make use of \(a\) ).
\item
  Let \(g_1, \ldots, g_p\) be in \(\mathcal{K}(\mathbf{R})\) . Deduce from \(b\) that there exist \(\mu \in \mathbb{M}\) and \(\nu \in \mathcal{M}(\mathrm{I})\) such that \(\nu(g_1) = \ldots = \nu(g_p) = 0\) , \((\nu * \mu)(f) \neq 0\) .
\item
  Deduce from c) that the mapping \((\mu,\nu)\mapsto\nu*\mu\) of \(\mathcal{M}(\mathbf{T})\times\mathcal{M}(\mathbf{T})\) into \(\mathcal{M}(\mathbf{T})\) is not vaguely continuous.
\end{enumerate}

\begin{enumerate}
\def\labelenumi{\arabic{enumi})}
\setcounter{enumi}{6}
\tightlist
\item
  Let \(G\) be a locally compact group.
\end{enumerate}

\begin{enumerate}
\def\labelenumi{\alph{enumi})}
\item
  Show that if a positive measure \(\mu \in \mathcal{C}'(\mathrm{G})\) admits an inverse in \(\mathcal{C}'(\mathrm{G})\) that is a positive measure, then \(\mu\) is necessarily a point measure.
\item
  Taking G to be the finite group \(\mathbf{Z}/2\mathbf{Z}\) , give an example of a non point measure \(\mu \in \mathcal{M}(\mathbf{G})\) that is invertible in \(\mathcal{M}(\mathbf{G})\) .
\end{enumerate}

\begin{enumerate}
\def\labelenumi{\arabic{enumi})}
\setcounter{enumi}{7}
\item
  Let G be a locally compact group; denote by \(\mathcal{C}_{+}^{\prime}(G)\) the set of positive measures on G with compact support. Show that the mapping \((\mu, \nu) \mapsto \mu * \nu\) of \(\mathcal{M}_{+}(G) \times \mathcal{C}_{+}^{\prime}(G)\) into \(\mathcal{M}_{+}(G)\) is continuous when \(\mathcal{C}'(G)\) is equipped with the weak topology \(\sigma(\mathcal{C}'(G), \mathcal{C}(G))\) , and \(\mathcal{M}(G)\) with the vague topology \(\sigma(\mathcal{M}(G), \mathcal{K}(G))\) . (Make use of Exer. 10 of Ch. III, §1 and the fact that G is paracompact.)
\item
  \begin{enumerate}
  \def\labelenumii{\alph{enumii})}
  \tightlist
  \item
    Let G be a locally compact group, B a bounded subset of \(\mathcal{M}(\mathbf{G})\) , C an equicontinuous subset of \(\mathcal{C}'(\mathbf{G})\) ; show that if \(\mathcal{C}'(\mathbf{G})\) is equipped with the weak topology \(\sigma(\mathcal{C}'(\mathbf{G}), \mathcal{C}(\mathbf{G}))\) , and \(\mathcal{M}(\mathbf{G})\) with the vague topology \(\sigma(\mathcal{M}(\mathbf{G}), \mathcal{K}(\mathbf{G}))\) , then the mapping \((\mu, \nu) \mapsto \mu * \nu\) of B × C into \(\mathcal{M}(\mathbf{G})\) is continuous. (Observe that all of the measures \(\nu \in \mathbf{C}\) have their support in a common compact set, and make use of Ch. III, §4, No.~3, remarks following Prop. 6.)
  \end{enumerate}
\end{enumerate}

\begin{enumerate}
\def\labelenumi{\alph{enumi})}
\setcounter{enumi}{1}
\tightlist
\item
  If, in \(\mathcal{M}^1 (\mathbf{R})\) , one sets \(\mu_n = \varepsilon_n\) , \(\nu_{n} = \varepsilon_{-n}\) , the sequences \((\mu_n)\) and \((\nu_{n})\) tend vaguely to 0 for the weak topology \(\sigma (\mathcal{M}^1 (\mathbf{R}),\overline{\mathcal{K}(\mathbf{R})})\) , but the sequence \((\mu_n*\nu_n)\) does not tend to 0 for the vague topology.
\end{enumerate}

¶ 10) For a locally compact space T, one denotes by \(\mathcal{C}^{\infty}(T)\) the Banach space of continuous bounded numerical functions on T. A subset H of \(\mathcal{M}^{1}(T)\) is said to be cramped if, for every \(\varepsilon > 0\) , there exists a compact subset K of T such that \(|\mu|(T - K) \leqslant \varepsilon\) for all \(\mu \in H\) .

\begin{enumerate}
\def\labelenumi{\alph{enumi})}
\item
  Show that if a subset H of \(\mathcal{M}^{1}(G)\) is cramped and is bounded for the topology defined by the norm of \(\mathcal{M}^{1}(T)\) , then it is relatively compact for the topology \(\sigma(\mathcal{M}^{1}(T),\mathcal{C}^{\infty}(T))\) (observe that H is relatively compact for \(\sigma(\mathcal{M}^{1}(T),\mathcal{K}(T))\) ).
\item
  Assuming in addition that T is paracompact, show that, conversely, if H is a subset of \(\mathcal{M}^1 (\mathrm{T})\) relatively compact for \(\sigma (\mathcal{M}^1 (\mathrm{T}),\mathcal{C}^\infty (\mathrm{T}))\) , then H is bounded for the norm of \(\mathcal{M}^1 (\mathrm{T})\) and is cramped. (Consider first the case that \(\mathbf{T} = \mathbf{N}\) , applying Exer. 15 of Ch. V, §5. Then contemplate the case that T is countable at infinity, the union of a sequence \((\mathrm{U}_n)\) of relatively compact open sets such that \(\overline{\mathrm{U}}_n\subset \mathrm{U}_{n + 1}\) . Arguing by contradiction, show that one can reduce to the case that there would exist for each \(n\) a continuous numerical function \(f_{n}\) defined on T, with support contained in \(\mathrm{U}_{n + 1} - \overline{\mathrm{U}}_n\) , such that \(\| f_n\| \leqslant 1\) , and a sequence \((\mu_n)\) of measures belonging to H and for which \(\mu_{n}(f_{n})\geqslant a > 0\) for all \(n\) . Consider then the continuous mapping \(u:\mathrm{L}^{1}(\mathbf{N})\to \mathcal{C}^{\infty}(\mathrm{T})\)
\end{enumerate}

such that \(u((\xi_n)) = \sum_{n=0}^{\infty} \xi_n f_n\) and obtain a contradiction to what was proved for \(T = N\) .

by considering the transpose of \(u\) . Finally, when \(T\) is an arbitrary paracompact locally compact space, \(T\) is the topological sum of a family \((T_{\alpha})\) of locally compact spaces that are countable at infinity; for every \(\alpha\) , let \(m_{\alpha} = \sup_{\mu \in H} |\mu|(T_{\alpha})\) ; show, arguing by contradiction and making use of the preceding case, that necessarily \(m_{\alpha} = 0\) except for a countable infinity of indices \(\alpha\) .)

\begin{enumerate}
\def\labelenumi{\alph{enumi})}
\setcounter{enumi}{2}
\tightlist
\item
  Show that the conclusion of \(b\) ) is not valid for the non-paracompact locally compact space defined in Exer. 16 h) of Ch. IV, §4.
\end{enumerate}

¶ 11) Let G be a locally compact group; on \(\mathcal{M}^{1}(G)\) , let us denote by \(\mathcal{T}_{\mathrm{I}}\) the topology \(\sigma(\mathcal{M}^{1}(G), \overline{\mathcal{K}(G)})\) , by \(\mathcal{T}_{\mathrm{III}}\) the topology \(\sigma(\mathcal{M}^{1}(G), \mathcal{C}^{\infty}(G))\) , and let us write \(\mathcal{M}_{\mathrm{I}}\) , \(\mathcal{M}_{\mathrm{III}}\) for the space \(\mathcal{M}^{1}(G)\) equipped respectively with \(\mathcal{T}_{\mathrm{I}}\) and \(\mathcal{T}_{\mathrm{III}}\) .

\begin{enumerate}
\def\labelenumi{\alph{enumi})}
\item
  Let A be a bounded subset of \(\mathcal{M}_{\mathrm{I}}\) , and B a relatively compact subset of \(\mathcal{M}_{\mathrm{III}}\) . Show that the restriction to \(\mathrm{A} \times \mathrm{B}\) of the mapping \((\mu, \nu) \mapsto \mu * \nu\) of \(\mathcal{M}_{\mathrm{I}} \times \mathcal{M}_{\mathrm{III}}\) into \(\mathcal{M}_{\mathrm{I}}\) is continuous (use Exer. 10 to reduce to evaluating an integral \(\iint f(st) d\mu(s) d\nu(t)\) when \(f(st) = \sum_{i} u_{i}(s)v_{i}(t)\) with \(u_{i}\) and \(v_{i}\) in \(\mathcal{K}(\mathrm{G})\) ).
\item
  Give an example where G is compact (hence \(T_{I} = T_{III}\) ) showing that \((\mu, \nu) \mapsto \mu * \nu\) , regarded as a mapping of \(M_{I} \times M_{III}\) into \(M_{I}\) , is not hypocontinuous with respect to the relatively compact subsets of \(M_{I}\) , nor with respect to the relatively compact subsets of \(M_{III}\) (cf.~Exer. 6).
\item
  Let \(A, B\) be two relatively compact subsets of \(\mathcal{M}_{\mathrm{III}}\) . Show that the restriction to \(A \times B\) of the mapping \((\mu, \nu) \mapsto \mu * \nu\) of \(\mathcal{M}_{\mathrm{III}} \times \mathcal{M}_{\mathrm{III}}\) into \(\mathcal{M}_{\mathrm{III}}\) is continuous (same method as in a)).
\item
  Denote by E the subspace of \(\mathbf{R}^{\mathrm{G}}\) formed by the linear combinations of characteristic functions of open subsets of G, by \(\mathcal{T}_{\mathrm{IV}}\) the topology \(\sigma(\mathcal{M}^{1}(\mathrm{G}),\mathrm{E})\) , and by \(\mathcal{M}_{\mathrm{IV}}\) the space \(\mathcal{M}^{1}(\mathrm{G})\) equipped with \(\mathcal{T}_{\mathrm{IV}}\) . Recall that the topologies induced by \(\mathcal{T}_{\mathrm{III}}\) and \(\mathcal{T}_{\mathrm{IV}}\) on a bounded subset of \(\mathcal{M}^{1}(\mathrm{G})\) consisting of positive measures, are in general distinct (Ch. V, §5, Exer. 16 c)). Recall also that the compact subsets of \(\mathcal{M}_{\mathrm{IV}}\) are the same as those of \(\mathcal{M}^{1}(\mathrm{G})\) for the weakened topology \(\sigma(\mathcal{M}^{1}(\mathrm{G}),(\mathcal{M}^{1}(\mathrm{G}))'\) on the Banach space \(\mathcal{M}^{1}(\mathrm{G})\) (Ch. VI, §2, Exer. 12). Show that if A, B are two relatively compact subsets of \(\mathcal{M}_{\mathrm{IV}}\) , the restriction to A × B of the mapping \((\mu,\nu)\mapsto\mu*\nu\) of \(\mathcal{M}_{\mathrm{IV}}\times\mathcal{M}_{\mathrm{IV}}\) into \(\mathcal{M}_{\mathrm{IV}}\) is continuous. (Restrict attention to the case that A and B are compact, and begin by proving that the image of A × B under the preceding mapping is then compact for \(\sigma(\mathcal{M}^{1}(\mathrm{G}),(\mathcal{M}^{1}(\mathrm{G}))'\) ); for this, apply the theorems of Eberlein and Šmulian (TVS, IV, §5, No.~3), and thus reduce to proving that if \((\mu_{n})\) , \((\nu_{n})\) are two sequences that converge to 0 in \(\mathcal{M}_{\mathrm{IV}}\) , then the same is true of the sequence \((\mu_{n}*\nu_{n})\) ; make use of Prop. 12 and Exer. 15 of Ch. V, §5. Finally, to prove that \((\mu,\nu)\mapsto\mu*\nu\) is continuous for \(\mathcal{T}_{\mathrm{IV}}\) , use c) and the fact that \(\mathcal{T}_{\mathrm{III}}\) is coarser than \(\mathcal{T}_{\mathrm{IV}}\) .)
\item
  Take \(\mathbf{G} = \mathbf{R}^2\) ; let \(a, b\) be the vectors of the canonical basis of \(\mathbf{G}\) over \(\mathbf{R}\) , \(\rho_n\) the measure on the interval \(\mathbf{I} = [0, \pi]\) of \(\mathbf{R}\) having as density with respect to Lebesgue measure the function \(\sin(2^n x)\) , and \(\mu_n\) the measure \(\rho_n \otimes \varepsilon_0\) on \(\mathbf{R} \times \mathbf{R}\) ; on the other hand let \(\nu_n = \varepsilon_{b/2^n} - \varepsilon_0\) on \(\mathbf{G}\) . Show that the sequence \((\mu_n)\) tends to 0 in \(\mathcal{M}_{\mathrm{IV}}\) and that the sequence \((\nu_n)\) tends to 0 in \(\mathcal{M}_{\mathrm{III}}\) , but that the sequence \((\mu_n * \nu_n)\) does not tend to 0 in \(\mathcal{M}_{\mathrm{IV}}\) (cf.~Ch. V, §5, Exer. 16).
\item
  Give an example where G is compact and where \((\mu,\nu)\mapsto\mu*\nu\) , regarded as a mapping of \(M_{IV}\times M_{IV}\) into \(M_{IV}\) , is not hypocontinuous relative to the compact subsets of \(M_{IV}\) (same method as in Exer. 6, on observing that for a given \(f\in E\) there exist compact subsets \(H\subset M_{IV}\) such that the set of functions \(t\mapsto\int f(st)d\mu(s)\) , where \(\mu\) runs over H, has finite rank over R).
\end{enumerate}

\begin{enumerate}
\def\labelenumi{\arabic{enumi})}
\setcounter{enumi}{11}
\tightlist
\item
  Let \(G\) be a locally compact group that is not unimodular.
\end{enumerate}

\begin{enumerate}
\def\labelenumi{\alph{enumi})}
\item
  Show that there exists a bounded positive measure \(\mu\) on \(\mathbf{G}\) such that \(\Delta_{\mathbf{G}} \cdot \mu\) is not bounded (take \(\mu\) to be discrete).
\item
  Let \(\mu'\) be a left Haar measure on G. Then \(\mu\) and \(\mu'\) are convolvable (Prop. 5). Show that \(\mu'\) and \(\mu\) are not convolvable.
\end{enumerate}

\begin{enumerate}
\def\labelenumi{\arabic{enumi})}
\setcounter{enumi}{12}
\tightlist
\item
  Let \(r\) be a number such that \(0 < r < 1\) ; for every integer \(n \geqslant 1\) , denote by \(\lambda_{n,r}\) the measure \((\varepsilon_{r^n} + \varepsilon_{-r^n}) / 2\) on \(\mathbf{R}\) , and set \(\mu_{n,r} = \lambda_{1,r} * \lambda_{2,r} * \dots * \lambda_{n,r}\) .
\end{enumerate}

\begin{enumerate}
\def\labelenumi{\alph{enumi})}
\item
  Show that the sequence \((\mu_{n,r})\) converges vaguely to a measure \(\mu_{r}\) on R, with support contained in I = {[}-1,+1{]} (prove that for every interval U of R, the sequence \((\mu_{n,r}(U))\) is convergent).
\item
  Show that for \(r < 1/2\) , the measure \(\mu_r\) is alien to the Lebesgue measure on \(\mathbf{R}\) , but that \(\mu_{1/2}\) is the measure induced on I by Lebesgue measure.
\item
  Let \(\nu_{1/4}\) be the image of \(\mu_{1/4}\) under the homothety \(t \mapsto 2t\) in \(\mathbf{R}\) . Show that \(\mu_{1/4} * \nu_{1/4} = \mu_{1/2}\) , even though \(\mu_{1/4}\) and \(\nu_{1/4}\) are both alien to Lebesgue measure (use Exer. 11 a)).
\end{enumerate}

